% Overleaf: Türkçe CV — main_tr.tex olarak yükle veya main.tex yap
% pdfLaTeX | ATS uyumlu
\documentclass[11pt,a4paper]{article}

\usepackage[utf8]{inputenc}
\usepackage[T1]{fontenc}
\usepackage[turkish]{babel}
\usepackage[margin=1.6cm]{geometry}
\usepackage{enumitem}
\usepackage{hyperref}

\hypersetup{colorlinks=true, urlcolor=black, linkcolor=black}
\pagestyle{empty}
\setlength{\parindent}{0pt}
\setlength{\parskip}{0pt}

\usepackage[scaled=0.95]{helvet}
\renewcommand{\familydefault}{\sfdefault}

\newcommand{\cvsection}[1]{%
  \vspace{0.7em}%
  {\normalsize\fontseries{m}\selectfont #1}\\[0.25em]
}
\newcommand{\jobentry}[4]{%
  \textbf{#1} \hfill {\small #2}\\
  {\small #3} \hfill {\small #4}\\[0.2em]
}

\begin{document}

{\LARGE Ece Dalpolat} \hfill Yapay Zeka Mühendisi \& Veri Bilimci\\[0.35em]
{\small
+90 552 360 9850 \,|\, ecedlplt9850@gmail.com \,|\, İstanbul \,|\,
\href{https://ecedalpolat.github.io/ece-dalpolat-portfolio/}{Portföy} \,|\,
\href{https://github.com/EceDalpolat}{GitHub} \,|\,
\href{https://www.linkedin.com/in/ece-dalpolat-0265a0221/}{LinkedIn}
}

\cvsection{Özet}

İK ve yetenek analitiği için üretim ortamında çalışan yapay zeka sistemleri kuruyorum; demo değil, canlı B2B ürünler. Qkare'de kurumsal müşterilerin kullandığı sistemlerin ana mühendisiyim: dashboard besleyen dbt veri katmanı, psikometrik veri üzerinde soru yanıtlayan LangChain ve LangGraph ajanları, KPI ve yetkinlik öneren doğrulanmış JSON çıktılı FastAPI servisleri.

Günlük işim: SQL'de veri modelleme (109+ dbt modeli, native PostgreSQL satır düzeyi güvenlik, Superset), LLM hatları tasarlama (RAG, semantik önbellek, prompt injection koruması), ortak agent gateway arkasında mikroservis yayınlama, canlıya almadan önce LLM-as-judge ile kalite ölçme. Belirsiz ihtiyaçtan Docker/Kubernetes ve CI ile çok kiracılı, maliyet kontrollü üretime kadar uçtan uca sahipleniyorum.

Işık Üniversitesi Bilgi Teknolojileri yüksek lisans öğrencisiyim. Miuul'da veri bilimi ve makine öğrenmesi dersleri veriyorum. Üretime çıkan yapay zeka mühendisliği, veri mühendisliği ve ML platform rollerine açığım.

\cvsection{Deneyim}

\jobentry{Yapay Zeka Mühendisi}{Haz 2025 -- Devam}{Qkare}{İstanbul}
\begin{itemize}[leftmargin=1.2em, nosep, topsep=2pt]
  \item Ana mühendis: dbt katmanı, üretim AI ajanları/danışmanları, gateway ve güvenlik; 8+ canlı sistem.
\end{itemize}

\jobentry{Öğretim Asistanı}{Oca 2025 -- Devam}{Miuul}{Uzaktan}
\begin{itemize}[leftmargin=1.2em, nosep, topsep=2pt]
  \item Veri bilimi ve makine öğrenmesi dersleri; uygulamalı proje mentörlüğü ve değerlendirme.
\end{itemize}

\jobentry{Lisansüstü Araştırma Görevlisi}{Oca 2024 -- Haz 2025}{Işık Üniversitesi}{İstanbul}
\begin{itemize}[leftmargin=1.2em, nosep, topsep=2pt]
  \item Yazılım laboratuvarları; veri bilimi ve derin öğrenme projelerinde danışmanlık.
\end{itemize}

\jobentry{Araştırma Veri Bilimci}{Oca 2024 -- Haz 2025}{Top4Honeychains}{}
\begin{itemize}[leftmargin=1.2em, nosep, topsep=2pt]
  \item Bal değer zinciri ETL, TensorFlow/scikit-learn modelleri, REST API entegrasyonları.
\end{itemize}

\cvsection{Projeler}

\textbf{Analitik Platformu} \hfill Kas 2025 -- May 2026\\
Qkare --- Uçtan uca İK analitiği: PostgreSQL FDW, 109 dbt modeli, 56 mart, 13 AI ön hesaplama tablosu, native RLS, Superset. Sorumluluk: dbt katmanı.

\vspace{0.25em}
\textbf{Analyzer Platformu} \hfill Nis 2025 -- May 2026\\
Qkare --- Psikometrik mart'lar üzerinde LangChain ReAct ajanı; 4 katman güvenlik, JWT+RLS, Langfuse. Sorumluluk: güvenlik ve config bootstrap.

\vspace{0.25em}
\textbf{Coach Platformu} \hfill Eyl 2025 -- Kas 2025\\
Qkare --- Çok ajanlı koçluk (LangGraph, POML, OpenAI Realtime). Referee Agent (LLM-as-judge), ses katmanı, eval araçları.

\vspace{0.25em}
\textbf{Skills Advisor} \hfill Ağu 2025 -- May 2026\\
Qkare --- Kapalı sözlükten tam 12 yetkinlik boyutu; LangChain veya DSPy ensemble. Ana mühendis.

\vspace{0.25em}
\textbf{KPI Advisor} \hfill Nis 2026 -- May 2026\\
Qkare --- APQC PCF ile SMART KPI; lexical retrieval, Golden Dataset, LLM-as-judge hattı. Ana mühendis.

\vspace{0.25em}
\textbf{Agent API Gateway} \hfill Nis 2026 -- May 2026\\
Qkare --- FastAPI monorepo, birleşik invoke API, JWT, ajan başına Docker/CI. Platform sorumlusu.

\vspace{0.25em}
\textbf{QKare AI Asistan} \hfill Tem 2025\\
Qkare --- Kurumsal chatbot: kural, semantik önbellek (FAISS/Milvus), RAG, LLM; hibrit güvenlik. Tek yazar.

\vspace{0.25em}
\textbf{Davranışsal Kümeleme Ar-Ge} \hfill Ağu 2025\\
Qkare --- K-Means, DBSCAN, GMM; PCA/t-SNE; Analitik Platform mart'larına girdi. Tek yazar.

\vspace{0.25em}
\textbf{CVAnalyzerAI} \hfill 2025\\
MCP ile özgeçmiş analizi (FastAPI, DSPy, ANTLR). github.com/EceDalpolat/cvAnalizerAI

\vspace{0.25em}
\textbf{TOP4HONEYCHAINS} \hfill 2024--2025\\
Bal izlenebilirlik ML, REST API, zaman serisi tahmini.

\vspace{0.25em}
\textbf{Pazarlama Kampanyası Tahmini}\\
Banka kampanya modeli: SMOTE, XGBoost, Streamlit.

\vspace{0.25em}
\textbf{Polen Sınıflandırma}\\
CNN, \%90+ doğruluk; veri artırma.

\vspace{0.25em}
\textbf{Anten Davranış Analizi}\\
DeepLabCut ile keypoint takibi; hareket analizi.

\vspace{0.25em}
\textbf{Öğrenci Başarı Analizi}\\
Akademik özelliklerle Random Forest; Tkinter arayüzü.

\cvsection{Beceriler}

\textbf{AI/LLM:} LangChain, LangGraph, DSPy, RAG, OpenAI API, Pydantic, LLM-as-judge\\
\textbf{Veri:} dbt Core, PostgreSQL, Apache Superset, Pandas, SQL\\
\textbf{Backend:} FastAPI, httpx, WebSocket, JWT, MCP, REST\\
\textbf{Altyapı:} Docker, Kubernetes, GitLab CI/CD, Redis, OpenTelemetry, Langfuse\\
\textbf{Kod:} Python, SQL, JavaScript, Java, C\#

\cvsection{Eğitim}

\textbf{Yüksek Lisans, Bilgi Teknolojileri} \hfill 2022 -- Devam\\
Işık Üniversitesi --- GNO 3,57. Tez: arıcılık ve bal hasadında makine öğrenmesi otomasyonu.

\vspace{0.25em}
\textbf{Lisans, Yazılım Mühendisliği} \hfill 2018 -- 2023\\
Kırklareli Üniversitesi --- GNO 2,71

\cvsection{Sertifikalar}

Miuul Veri Bilimci Bootcamp; Miuul Makine Öğrenmesi; Miuul Üretken AI; Udemy LangChain/RAG; BTK Uygulamalı SQL; Huawei HCCDA-AI

\vspace{0.3em}
{\small Türkçe (ana dil) | İngilizce (B2)}

\end{document}
